% =====================================================================
%  Explicacion detallada de la interfaz MATLAB (braille_gui)
%  Proyecto: Sistema bidireccional Texto <-> Braille en FPGA (DE10-Lite)
%  Compilar con:  pdflatex matlab_explicado.tex   (dos veces, para el indice)
% =====================================================================
\documentclass[11pt,a4paper]{article}

\usepackage[utf8]{inputenc}
\usepackage[T1]{fontenc}
\usepackage[spanish,es-noshorthands,es-tabla]{babel}
\usepackage{lmodern}
\usepackage{amssymb}
\usepackage[margin=2.3cm]{geometry}
\usepackage{xcolor}
\usepackage{listings}
\usepackage{booktabs}
\usepackage{tabularx}
\usepackage{longtable}
\usepackage{array}
\usepackage{enumitem}
\usepackage[most]{tcolorbox}
\usepackage{fancyhdr}
\usepackage{titlesec}
\usepackage[hidelinks]{hyperref}

% ---------------- Colores ----------------
\definecolor{tinta}{HTML}{13202B}
\definecolor{verde}{HTML}{137366}
\definecolor{verdesuave}{HTML}{DCEFEB}
\definecolor{ambar}{HTML}{9A520C}
\definecolor{ambarsuave}{HTML}{F8E6D3}
\definecolor{gris}{HTML}{5B6874}
\definecolor{fondocodigo}{HTML}{F4F6F7}
\definecolor{comentario}{HTML}{2F7D32}
\definecolor{cadena}{HTML}{A0306A}
\definecolor{clave}{HTML}{1F4E9A}

% ---------------- Titulos ----------------
\titleformat{\section}{\Large\bfseries\color{tinta}}{\thesection}{0.8em}{}[\vspace{-0.4em}\textcolor{verde}{\rule{\linewidth}{1pt}}]
\titleformat{\subsection}{\large\bfseries\color{tinta}}{\thesubsection}{0.7em}{}
\titleformat{\subsubsection}{\normalsize\bfseries\color{verde}}{\thesubsubsection}{0.6em}{}

% ---------------- Encabezado ----------------
\pagestyle{fancy}
\fancyhf{}
\fancyhead[L]{\small\color{gris}Interfaz MATLAB \textendash{} Sistema Braille FPGA}
\fancyhead[R]{\small\color{gris}\leftmark}
\fancyfoot[C]{\small\thepage}
\renewcommand{\headrulewidth}{0.4pt}

% ---------------- Codigo MATLAB ----------------
\lstdefinestyle{matlab}{
  language=Matlab,
  basicstyle=\ttfamily\small,
  keywordstyle=\color{clave}\bfseries,
  commentstyle=\color{comentario}\itshape,
  stringstyle=\color{cadena},
  backgroundcolor=\color{fondocodigo},
  frame=single, rulecolor=\color{gris!40},
  framesep=6pt, xleftmargin=8pt, xrightmargin=8pt,
  numbers=left, numberstyle=\tiny\color{gris}, numbersep=6pt,
  breaklines=true, breakatwhitespace=true,
  showstringspaces=false, columns=fullflexible, keepspaces=true,
  morekeywords={uifigure,uigridlayout,uilabel,uidropdown,uibutton,uicheckbox,uipanel,uieditfield,uitextarea,uispinner,uiaxes,serialport,serialportlist,configureTerminator,configureCallback,flush,write,read,uialert,hex2dec,sprintf,isempty,numel,find,strtrim,datetime},
  literate={á}{{\'a}}1 {é}{{\'e}}1 {í}{{\'i}}1 {ó}{{\'o}}1 {ú}{{\'u}}1 {ñ}{{\~n}}1 {→}{{$\rightarrow$}}1
}
\lstset{style=matlab}
\newcommand{\linea}[1]{\textcolor{gris}{\small(líneas #1)}}

% ---------------- Cajas ----------------
\newtcolorbox{idea}[1][Idea clave]{colback=verdesuave,colframe=verde,fonttitle=\bfseries,title=#1,arc=2mm,boxrule=0.8pt}
\newtcolorbox{ojo}[1][Atención]{colback=ambarsuave,colframe=ambar,fonttitle=\bfseries,title=#1,arc=2mm,boxrule=0.8pt}
\newtcolorbox{pregunta}[1]{colback=white,colframe=tinta,fonttitle=\bfseries,title={P: #1},arc=2mm,boxrule=0.6pt,breakable}
\newtcolorbox{diagrama}{fontupper=\footnotesize,colback=fondocodigo,colframe=gris!40,arc=2mm,boxrule=0.6pt,left=6pt,right=6pt}

\newcommand{\cod}[1]{\texttt{#1}}
\setlist{itemsep=2pt,topsep=4pt}
\setlength{\emergencystretch}{3em}
\renewcommand{\arraystretch}{1.25}

% =====================================================================
\begin{document}

% ---------------- Portada ----------------
\begin{titlepage}
\centering
\vspace*{3cm}
{\color{verde}\large\bfseries PROYECTO FINAL \textendash{} DE10-LITE\par}
\vspace{0.8cm}
{\color{tinta}\Huge\bfseries Interfaz MATLAB del\\[0.2cm] Sistema Braille FPGA\par}
\vspace{0.6cm}
{\Large\color{gris} Explicación completa de \texttt{braille\_gui}\par}
\vspace{1.5cm}
\begin{minipage}{0.8\textwidth}
\begin{idea}[Qué contiene este documento]
Qué hace el programa, los conceptos de MATLAB que usa, cada variable, cada
elemento de la ventana, cada función línea por línea, el protocolo de
comunicación con la FPGA, los flujos completos de cada modo y una sección de
preguntas frecuentes para responder en una presentación.
\end{idea}
\end{minipage}
\vfill
{\color{gris} Archivo del proyecto: \texttt{matlab} (se ejecuta como \texttt{braille\_gui.m})\par}
\end{titlepage}

\tableofcontents
\newpage

% =====================================================================
\section{Panorama general}

\subsection{¿Qué es \texttt{braille\_gui}?}
Es la \textbf{ventana de la computadora} que se comunica con la FPGA por el
puerto serial (UART). Está escrita en MATLAB usando \emph{App Designer
programático} (componentes \cod{ui...}) y no necesita toolboxes adicionales.

\begin{idea}
Toda la lógica del sistema vive en la FPGA. MATLAB \textbf{no decide nada}:
solo \textbf{escucha} lo que reporta la FPGA, lo \textbf{muestra} en pantalla,
lo \textbf{dice en voz alta} y, en dos casos, \textbf{envía} datos.
\end{idea}

\subsection{Qué hace y qué no hace}
\begin{center}
\begin{tabularx}{\linewidth}{>{\raggedright}X >{\raggedright\arraybackslash}X}
\toprule
\textbf{MATLAB sí hace} & \textbf{MATLAB no hace (lo hace la FPGA)} \\
\midrule
Abrir el puerto serial a 9600 baudios & Elegir el modo (lo eligen los switches) \\
Mostrar el panel del modo actual & Convertir letras en puntos Braille \\
Enviar la palabra en modo lectura (\cod{R}) & Elegir la letra de práctica (LFSR) \\
Enviar cuántas letras tiene una sesión (\cod{L}) & Comparar y contar aciertos y errores \\
Preguntar el modo al conectar (\cod{Q}) & Decidir cuándo termina una sesión \\
Calcular el porcentaje y dibujar la gráfica & Mover los servos y leer los botones \\
Hablar con la voz de Windows & \\
\bottomrule
\end{tabularx}
\end{center}

\subsection{Dónde está dentro del sistema}
\begin{diagrama}
\begin{verbatim}
   +-------------------+      USB / Bluetooth      +------------------------+
   |   PC (MATLAB)     |  <--------------------->  |   FPGA DE10-Lite       |
   |   braille_gui     |     UART 9600 8N1         |   braille_top          |
   |                   |                           |                        |
   |  envia:  Q, R, L  |  ---- texto + \n ---->    |  line_rx -> controller |
   |  recibe: M D E K  |  <--- texto + \n -----    |  controller -> msg_tx  |
   |          T O X F  |                           |                        |
   +-------------------+                           +------------------------+
\end{verbatim}
\end{diagrama}

\subsection{Requisitos para ejecutarlo}
\begin{itemize}
  \item \textbf{MATLAB R2019b o más nuevo}: \cod{serialport} y los componentes
        \cod{ui...} aparecieron en esa versión.
  \item \textbf{Windows} para la voz (usa la librería \cod{System.Speech} de .NET).
        En otro sistema el programa funciona, solo que sin voz.
  \item \textbf{El nombre del archivo}: en el repositorio se llama \cod{matlab}, sin
        extensión. Para ejecutarlo debe guardarse como \cod{braille\_gui.m}, porque
        MATLAB exige que el archivo se llame igual que su función principal.
  \item Se ejecuta escribiendo \cod{braille\_gui} en la ventana de comandos.
\end{itemize}

% =====================================================================
\section{Conceptos de MATLAB que usa el programa}

\subsection{Sintaxis básica}
\begin{longtable}{>{\ttfamily}p{3.6cm} p{10.6cm}}
\toprule
\normalfont\textbf{Escritura} & \textbf{Significado} \\
\midrule
\endhead
\% texto & Comentario: MATLAB lo ignora. \\
x = 5; & Asigna un valor. El \cod{;} evita que se imprima en la consola. \\
'hola' & Texto tipo \textbf{char} (comillas simples): un arreglo de caracteres. \\
"hola" & Texto tipo \textbf{string} (comillas dobles). Se pueden unir con \cod{+}: \cod{"R" + "hola"} $\rightarrow$ \cod{"Rhola"}. \\
{[}a b{]} & Une elementos en un arreglo. Con char: \cod{['ho' 'la']} $\rightarrow$ \cod{'hola'}. \\
m(2) & Elemento número 2. \textbf{En MATLAB los índices empiezan en 1}, no en 0. \\
m(3:4) & Elementos del 3 al 4. \\
end (en índice) & El último elemento: \cod{x(end+1) = v} agrega \cod{v} al final. \\
{[}{]} \ y \ '' & Vacío. \cod{isempty(x)} pregunta si está vacío. \\
uint8(x) & Convierte a entero de 0 a 255, o sea un \textbf{byte}. \cod{uint8('A')} $\rightarrow$ 65. \\
char(x) & Convierte números a caracteres. \cod{char(65)} $\rightarrow$ \cod{'A'}. \\
\textasciitilde & Negación lógica (\cod{\textasciitilde isempty(s)}) o ``ignorar argumento'' (\cod{f(\textasciitilde, \textasciitilde)}). \\
@nombre & Referencia a una función, para que otro la llame después (callback). \\
@(a,b) expr & Función anónima: una mini función escrita en una línea. \\
{[}r g b{]} & Color con valores de 0 a 1. \cod{[0.75 0 0]} es rojo oscuro. \\
\bottomrule
\end{longtable}

\subsection{Funciones anidadas: variables compartidas}
Todo el programa es \textbf{una función grande} (\cod{braille\_gui}) que contiene
\textbf{funciones pequeñas adentro} (\cod{conectar}, \cod{enviar}, \cod{procesar}, \ldots).

\begin{lstlisting}
function braille_gui()
    s = [];                       % variable de la funcion grande
    function conectar(~, ~)
        s = serialport(...);      % la funcion de adentro la modifica
    end
end
\end{lstlisting}

\begin{idea}
Las funciones anidadas \textbf{comparten las variables} de la función que las
contiene. Por eso \cod{s}, \cod{rxbuf}, \cod{texto}, \cod{histAcc}, etc. se pueden
leer y modificar desde cualquier parte del programa sin pasarlas como
argumento. Funcionan como ``variables globales'' pero solo dentro de este
programa.
\end{idea}

\subsection{Programación por eventos (callbacks)}
Un \textbf{callback} es una función que MATLAB ejecuta \textbf{automáticamente}
cuando ocurre algo:

\begin{center}
\begin{tabular}{>{\raggedright}p{6.4cm} >{\ttfamily\raggedright\arraybackslash}p{6.8cm}}
\toprule
\textbf{Evento} & \normalfont\textbf{Callback que se ejecuta} \\
\midrule
Presionar el botón Conectar & conectar \\
Presionar ``Enviar palabra'' & enviarPalabra \\
Cambiar el número de letras & cambiarN \\
Presionar ``Leer en voz alta'' & @(\textasciitilde,\textasciitilde) hablar(texto) \\
Presionar ``Borrar'' & borrarTexto \\
Llega un \textbackslash n por el puerto serial & alRecibir \\
Cerrar la ventana & cerrar \\
\bottomrule
\end{tabular}
\end{center}

\begin{idea}[Consecuencia importante]
El programa \textbf{no tiene un ciclo principal}. Después de construir la ventana,
la función \cod{braille\_gui} termina y MATLAB se queda esperando eventos. Cada
botón o cada dato que llega dispara su callback.
\end{idea}

% =====================================================================
\section{Mapa del archivo}

\begin{center}
\begin{tabular}{>{\raggedright}p{3.2cm} >{\raggedright}p{4.6cm} >{\raggedright\arraybackslash}p{6.2cm}}
\toprule
\textbf{Líneas} & \textbf{Bloque} & \textbf{Contenido} \\
\midrule
1--19 & Encabezado & Comentario con el propósito y el uso \\
21--32 & Configuración & Constantes y variables de estado \\
34--129 & Ventana & Construcción de toda la interfaz \\
135--186 & Conexión y recepción & \cod{listaPuertos}, \cod{conectar}, \cod{enviar}, \cod{alRecibir} \\
192--315 & Mensajes de la FPGA & \cod{procesar}, \cod{mostrarModo}, \cod{nombreModo} \\
321--332 & Práctica & \cod{enviarN}, \cod{cambiarN} \\
338--348 & Lectura & \cod{enviarPalabra} \\
354--376 & Escritura & \cod{letraEscrita}, \cod{borrarTexto} \\
382--409 & Utilidades & \cod{registrar}, \cod{hablar}, \cod{cerrar} \\
\bottomrule
\end{tabular}
\end{center}

% =====================================================================
\section{Configuración y variables de estado}
\linea{21--32}

\begin{lstlisting}[firstnumber=21]
    BAUD    = 9600;
    MAX_LEN = 16;          % Debe coincidir con MAX_LEN en la FPGA

    s      = [];           % serialport
    rxbuf  = uint8([]);    % bytes recibidos sin procesar
    synth  = [];           % voz
    texto  = '';           % escritura
    palabraW = '';
    histAcc   = [];           % precision tras cada intento (practica)
    modoAct   = 'I';          % modo que reporto la FPGA
    nLetras   = 10;           % letras por sesion (lo cuenta la FPGA)
    intentos  = 0;            % intentos hechos en la sesion
\end{lstlisting}

\begin{longtable}{>{\ttfamily}p{2.3cm} p{4.6cm} p{7.2cm}}
\toprule
\normalfont\textbf{Variable} & \textbf{Qué guarda} & \textbf{Por qué existe} \\
\midrule
\endhead
BAUD & 9600 & Velocidad del puerto. \textbf{Debe ser igual} a la de la FPGA (\cod{CLKS\_PER\_BIT = 50\,000\,000 / 9600}). \\
MAX\_LEN & 16 & Máximo de letras por palabra. Igual que en la FPGA, para avisar antes de enviar una palabra demasiado larga. \\
s & El objeto del puerto serial & Vacío (\cod{[]}) significa ``no conectado''. \\
rxbuf & Bytes recibidos sin procesar & Los mensajes pueden llegar en pedazos; aquí se acumulan hasta tener un \cod{\textbackslash n}. \\
synth & El sintetizador de voz & Se crea la primera vez que el programa habla y se reutiliza. \\
texto & Todo lo escrito en modo escritura & Para mostrarlo en el área de texto. \\
palabraW & La palabra actual en escritura & Para leerla completa en voz alta cuando llega un espacio. \\
histAcc & Lista de precisiones (\%) & Cada intento agrega un valor; con ella se dibuja la gráfica. \\
modoAct & \cod{'I'}, \cod{'R'}, \cod{'W'} o \cod{'P'} & El modo que reportó la FPGA. Se usa para saber si enviar \cod{L}. \\
nLetras & Letras por sesión & Para mostrar ``Letra 3 de 10''. \\
intentos & Intentos de la sesión & Ídem. Se calcula como aciertos + errores. \\
\bottomrule
\end{longtable}

% =====================================================================
\section{La ventana}
\linea{34--129}

\subsection{Distribución general}
\begin{diagrama}
\begin{verbatim}
+---------------------------- Braille FPGA ----------------------------+
| Puerto: [COM5 v] [refrescar] [Conectar]  Desconectado        [x] Voz | Fila 1 (34 px)
+----------------------------------------------------------------------+
|                         MODO: LECTURA                                | Fila 2 (50 px)
+----------------------------------------------------------------------+
|                                                                      |
|     Un solo panel visible segun el modo:                             | Fila 3 ('1x')
|     Reposo / Lectura / Escritura / Practica                          |
|                                                                      |
+----------------------------------------------------------------------+
| [12:01:05] Conectado a COM5.                                         | Fila 4 (110 px)
+----------------------------------------------------------------------+
\end{verbatim}
\end{diagrama}

\begin{lstlisting}[firstnumber=35]
    fig = uifigure('Name', 'Braille FPGA', 'Position', [100 80 900 640]);
    fig.CloseRequestFcn = @cerrar;

    g = uigridlayout(fig, [4 1]);
    g.RowHeight = {34, 50, '1x', 110};
\end{lstlisting}

\begin{itemize}
  \item \cod{uifigure}: crea la ventana con título ``Braille FPGA'', en la posición
        $(100, 80)$ y de $900 \times 640$ píxeles.
  \item \cod{CloseRequestFcn = @cerrar}: al cerrar la ventana se ejecuta
        \cod{cerrar}, que libera el puerto.
  \item \cod{uigridlayout(fig, [4 1])}: divide la ventana en una cuadrícula de 4
        filas y 1 columna. Los elementos se acomodan solos.
  \item \cod{RowHeight = \{34, 50, '1x', 110\}}: alto de cada fila. \cod{'1x'}
        significa ``todo el espacio que sobre''.
\end{itemize}

\subsection{Fila 1: barra de conexión}
\begin{lstlisting}[firstnumber=42]
    g1 = uigridlayout(g, [1 6]);
    g1.ColumnWidth = {55, 110, 34, 100, '1x', 60};
    g1.Padding = [0 0 0 0];
    uilabel(g1, 'Text', 'Puerto:');
    ddPort = uidropdown(g1, 'Items', listaPuertos());
    uibutton(g1, 'Text', char(8635), 'Tooltip', 'Actualizar puertos', ...
             'ButtonPushedFcn', @(~, ~) set(ddPort, 'Items', listaPuertos()));
    btnCon = uibutton(g1, 'Text', 'Conectar', 'ButtonPushedFcn', @conectar);
    lblEst = uilabel(g1, 'Text', 'Desconectado', 'FontColor', [0.75 0 0]);
    chkVoz = uicheckbox(g1, 'Text', 'Voz', 'Value', true);
\end{lstlisting}

\begin{center}
\begin{tabular}{>{\ttfamily}l l >{\raggedright\arraybackslash}p{7.5cm}}
\toprule
\normalfont\textbf{Variable} & \textbf{Componente} & \textbf{Función} \\
\midrule
--- & \cod{uilabel} & Texto fijo ``Puerto:'' \\
ddPort & \cod{uidropdown} & Lista desplegable para elegir el COM \\
--- & \cod{uibutton} & Botón $\circlearrowright$ (\cod{char(8635)}) que vuelve a buscar puertos \\
btnCon & \cod{uibutton} & Conectar / Desconectar (el mismo botón) \\
lblEst & \cod{uilabel} & Estado: rojo ``Desconectado'' o verde ``Conectado COM5'' \\
chkVoz & \cod{uicheckbox} & Activa o desactiva la voz \\
\bottomrule
\end{tabular}
\end{center}

El botón de refrescar usa una \textbf{función anónima}:
\cod{@(\textasciitilde, \textasciitilde) set(ddPort, 'Items', listaPuertos())}. Al
presionarlo, pide de nuevo la lista de puertos y la pone en el desplegable.

\subsection{Fila 2: modo actual}
\begin{lstlisting}[firstnumber=54]
    lblModo = uilabel(g, 'Text', 'MODO: -   (se elige con SW1-SW0 en la tarjeta)', ...
                      'FontSize', 24, 'FontWeight', 'bold', ...
                      'HorizontalAlignment', 'center');
\end{lstlisting}
Un texto grande y centrado. \cod{mostrarModo} lo cambia a ``MODO: LECTURA'', etc.

\subsection{Fila 3: un panel por modo}
\begin{idea}[El truco de los paneles]
Los cuatro paneles se colocan \textbf{en la misma celda}, uno encima de otro. Solo
se hace visible (\cod{Visible = true}) el del modo actual. Así la ventana cambia
de contenido sin crear ni destruir componentes.
\end{idea}

\begin{lstlisting}[firstnumber=59]
    area = uigridlayout(g, [1 1]);         % una sola celda
    pI = uipanel(area, 'BorderType', 'none');   % Reposo
    pR = uipanel(area, ...); pR.Layout.Row = 1; pR.Layout.Column = 1;  % Lectura
    pW = uipanel(area, ...); ...           % Escritura (misma celda)
    pP = uipanel(area, ...); ...           % Practica  (misma celda)
\end{lstlisting}

\subsubsection{Panel Reposo (\texttt{pI})}
Solo un texto con instrucciones: qué combinación de switches elige cada modo.

\subsubsection{Panel Lectura (\texttt{pR})}
\begin{center}
\begin{tabular}{>{\ttfamily}l l >{\raggedright\arraybackslash}p{8cm}}
\toprule
\normalfont\textbf{Variable} & \textbf{Componente} & \textbf{Función} \\
\midrule
edPal & \cod{uieditfield} & Campo donde se escribe la palabra (máximo 16) \\
--- & \cod{uibutton} & ``Enviar palabra'' $\rightarrow$ llama a \cod{enviarPalabra} \\
lblLetra & \cod{uilabel} & La letra que muestran los servos, en tamaño 150 \\
lblProg & \cod{uilabel} & Mensaje de ayuda o progreso \\
\bottomrule
\end{tabular}
\end{center}

\subsubsection{Panel Escritura (\texttt{pW})}
\begin{center}
\begin{tabular}{>{\ttfamily}l l >{\raggedright\arraybackslash}p{8cm}}
\toprule
\normalfont\textbf{Variable} & \textbf{Componente} & \textbf{Función} \\
\midrule
taTexto & \cod{uitextarea} & Todo lo tecleado en Braille (no editable) \\
lblUlt & \cod{uilabel} & Última letra recibida o ``Patrón no reconocido'' \\
--- & \cod{uibutton} & ``Leer en voz alta'' $\rightarrow$ \cod{hablar(texto)} \\
--- & \cod{uibutton} & ``Borrar'' $\rightarrow$ \cod{borrarTexto} \\
\bottomrule
\end{tabular}
\end{center}

\subsubsection{Panel Práctica (\texttt{pP})}
\begin{center}
\begin{tabular}{>{\ttfamily}l l >{\raggedright\arraybackslash}p{8cm}}
\toprule
\normalfont\textbf{Variable} & \textbf{Componente} & \textbf{Función} \\
\midrule
spnN & \cod{uispinner} & Letras por sesión, de 0 a 99 (0 = sin límite). Al cambiar llama a \cod{cambiarN} \\
lblRes & \cod{uilabel} & ``Letra 3 de 10'', ``Correcto!'', ``Incorrecto: escribiste K'' \\
lblObj & \cod{uilabel} & Letra objetivo en grande; verde si acierto, rojo si error \\
lblStats & \cod{uilabel} & Aciertos, errores y precisión \\
ax & \cod{uiaxes} & Gráfica de precisión acumulada (eje Y de 0 a 100 \%) \\
\bottomrule
\end{tabular}
\end{center}

\subsection{Fila 4: registro}
\begin{lstlisting}[firstnumber=126]
    taLog = uitextarea(g, 'Editable', 'off', 'FontName', 'Consolas', 'Value', {''});
\end{lstlisting}
Un área de texto con todo lo que ocurre, con hora. La llena \cod{registrar}.

\subsection{Al terminar de construir}
\begin{lstlisting}[firstnumber=128]
    mostrarModo('I');
    registrar('Selecciona el puerto y presiona Conectar.');
\end{lstlisting}
Muestra el panel de Reposo y escribe la primera línea del registro. Desde aquí,
el programa solo reacciona a eventos.

% =====================================================================
\section{Conexión y envío}

\subsection{\texttt{listaPuertos}: buscar puertos COM}
\linea{135--138}
\begin{lstlisting}[firstnumber=135]
    function lst = listaPuertos()
        lst = cellstr(serialportlist("available"));
        if isempty(lst), lst = {'COM5'}; end
    end
\end{lstlisting}
\begin{itemize}
  \item \cod{serialportlist("available")}: pide al sistema los puertos seriales libres.
  \item \cod{cellstr}: los convierte en el formato que necesita el desplegable.
  \item Si no hay ninguno, pone \cod{'COM5'} como sugerencia para que la lista no quede vacía.
\end{itemize}

\subsection{\texttt{conectar}: abrir o cerrar el puerto}
\linea{140--164}
\begin{lstlisting}[firstnumber=140]
    function conectar(~, ~)
        if ~isempty(s)                       % desconectar
            configureCallback(s, "off");
            delete(s); s = [];
            btnCon.Text = 'Conectar';
            lblEst.Text = 'Desconectado'; lblEst.FontColor = [0.75 0 0];
            registrar('Desconectado.');
            return;
        end
        try
            s = serialport(ddPort.Value, BAUD);
        catch err
            s = [];
            uialert(fig, err.message, 'No se pudo abrir el puerto');
            return;
        end
        configureTerminator(s, "LF");
        flush(s);
        rxbuf = uint8([]);
        configureCallback(s, "terminator", @alRecibir);
        btnCon.Text = 'Desconectar';
        lblEst.Text = ['Conectado ' ddPort.Value]; lblEst.FontColor = [0 0.5 0];
        registrar(['Conectado a ' ddPort.Value '.']);
        enviar("Q");                         % la FPGA responde con su modo
    end
\end{lstlisting}

\begin{longtable}{p{2.2cm} p{11.9cm}}
\toprule
\textbf{Paso} & \textbf{Qué hace} \\
\midrule
\endhead
A. Desconectar (141--148) & Si \cod{s} no está vacío, ya había conexión: apaga el callback, cierra el puerto (\cod{delete}), vacía \cod{s}, regresa el botón a ``Conectar'' y el estado a rojo. \cod{return} sale de la función. \\
B. Abrir (149--155) & \cod{serialport(puerto, 9600)} abre el COM elegido. Por defecto usa 8 bits de datos, sin paridad y 1 bit de stop (\textbf{8N1}), igual que la FPGA. \cod{try/catch}: si el puerto no existe o está ocupado, muestra una ventana de error (\cod{uialert}) en vez de detener el programa. \\
C. Configurar (156--159) & \cod{configureTerminator(s, "LF")}: el fin de mensaje es el byte 10 (\cod{\textbackslash n}), igual que en la FPGA. \cod{flush}: borra basura previa. \cod{configureCallback(s, "terminator", @alRecibir)}: \textbf{cada vez que llegue un \cod{\textbackslash n}, ejecutar \cod{alRecibir}}. \\
D. Interfaz (160--162) & El botón dice ``Desconectar'' y el estado se pone verde. \\
E. Preguntar (163) & Envía \cod{Q}. MATLAB no puede ver los switches de la tarjeta, así que le pregunta a la FPGA en qué modo está. \\
\bottomrule
\end{longtable}

\subsection{\texttt{enviar}: la única función que manda datos}
\linea{166--172}
\begin{lstlisting}[firstnumber=166]
    function enviar(txt)
        if isempty(s)
            registrar('No conectado.');
            return;
        end
        write(s, [uint8(char(txt)) 10], "uint8");
    end
\end{lstlisting}

\begin{center}
\begin{tabular}{>{\ttfamily}l >{\raggedright}p{6cm} >{\ttfamily\raggedright\arraybackslash}p{4.8cm}}
\toprule
\normalfont\textbf{Parte} & \textbf{Qué hace} & \normalfont\textbf{Con \cod{"Rhola"}} \\
\midrule
char(txt) & string $\rightarrow$ caracteres & 'Rhola' \\
uint8(...) & caracteres $\rightarrow$ bytes & [82 104 111 108 97] \\
{[}... 10{]} & agrega el byte 10 (\cod{\textbackslash n}) & [82 104 111 108 97 10] \\
write(s, ..., "uint8") & manda los bytes por el puerto & 6 bytes por el cable \\
\bottomrule
\end{tabular}
\end{center}

\begin{ojo}[¿Por qué \texttt{uint8(char(...))}?]
Así cada carácter ocupa \textbf{exactamente un byte}. La \textbf{ñ} queda como el
byte 241 (codificación Latin-1), que es el valor que la FPGA reconoce en
\cod{ascii\_to\_braille}. Si se enviara en UTF-8 ocuparía 2 bytes y la FPGA no la
reconocería.
\end{ojo}

% =====================================================================
\section{Recepción: \texttt{alRecibir}}
\linea{175--186}

\begin{lstlisting}[firstnumber=175]
    function alRecibir(src, ~)
        n = src.NumBytesAvailable;
        if n == 0, return; end
        rxbuf = [rxbuf reshape(read(src, n, "uint8"), 1, [])];
        while true
            k = find(rxbuf == 10, 1);
            if isempty(k), break; end
            linea = rxbuf(1:k-1);
            rxbuf(1:k) = [];
            if ~isempty(linea), procesar(linea); end
        end
    end
\end{lstlisting}

\begin{idea}
Esta función es el equivalente de \cod{line\_rx} de la FPGA, pero del lado de la
PC: junta bytes hasta encontrar un \cod{\textbackslash n} y entrega cada mensaje
completo a \cod{procesar}.
\end{idea}

\begin{center}
\begin{tabular}{l >{\raggedright\arraybackslash}p{12cm}}
\toprule
\textbf{Línea} & \textbf{Qué hace} \\
\midrule
176--177 & Cuántos bytes hay esperando. Si ninguno, sale. \\
178 & Los lee todos y los \textbf{agrega al final} de \cod{rxbuf}. \cod{reshape(..., 1, [])} los acomoda en una fila. \\
179 & Ciclo: puede haber \textbf{varios mensajes juntos} en el buffer. \\
180--181 & Busca el primer byte 10. \cod{k} es su posición. Si no hay, sale del ciclo. \\
182 & El mensaje es todo lo que está \textbf{antes} del \cod{\textbackslash n}. \\
183 & Borra del buffer el mensaje \textbf{y} su \cod{\textbackslash n}. \\
184 & Si el mensaje no está vacío, lo procesa. \\
\bottomrule
\end{tabular}
\end{center}

\subsection*{Ejemplo: llegan dos mensajes juntos}
\begin{diagrama}
\begin{verbatim}
rxbuf = [ 'D' 'h' 10 'D' 'o' 10 ]
vuelta 1: k = 3 -> linea = "Dh" -> procesar("Dh") -> rxbuf = [ 'D' 'o' 10 ]
vuelta 2: k = 3 -> linea = "Do" -> procesar("Do") -> rxbuf = [ ]
vuelta 3: no hay 10 -> break
\end{verbatim}
\end{diagrama}

\begin{ojo}[¿Por qué un buffer propio?]
Los bytes pueden llegar en pedazos: por ejemplo \cod{D} en una lectura y
\cod{h\textbackslash n} en la siguiente. El pedazo incompleto se queda en
\cod{rxbuf} hasta que llega su \cod{\textbackslash n}, así ningún mensaje se procesa
a medias.
\end{ojo}

% =====================================================================
\section{Procesamiento de mensajes: \texttt{procesar}}
\linea{192--286}

\subsection{Preparación}
\begin{lstlisting}[firstnumber=192]
    function procesar(m)
        c2 = ' ';
        if numel(m) >= 2, c2 = char(m(2)); end

        switch char(m(1))
\end{lstlisting}
\begin{itemize}
  \item \cod{m}: el mensaje en bytes, sin el \cod{\textbackslash n}.
  \item \cod{c2}: el \textbf{segundo carácter}, que casi siempre es el dato (la letra o
        el modo). Si el mensaje tiene un solo carácter (como \cod{E}), queda en espacio.
  \item \cod{switch char(m(1))}: decide según la \textbf{primera letra}, que indica el
        tipo de mensaje.
\end{itemize}

\subsection{Cómo se leen las posiciones}
\begin{center}
\begin{tabular}{l c c c c c c}
\toprule
\textbf{Posición en MATLAB} & 1 & 2 & 3 & 4 & 5 & 6 \\
\midrule
Mensaje \cod{Om0301} & O & m & 0 & 3 & 0 & 1 \\
Significado & tipo & letra & \multicolumn{2}{c}{aciertos} & \multicolumn{2}{c}{errores} \\
Se lee con & \cod{m(1)} & \cod{c2} & \multicolumn{2}{c}{\cod{m(3:4)}} & \multicolumn{2}{c}{\cod{m(5:6)}} \\
\bottomrule
\end{tabular}
\end{center}

\subsection{Tabla resumen de mensajes}
\begin{longtable}{>{\ttfamily}p{2cm} p{3.6cm} p{5.6cm} p{2.6cm}}
\toprule
\normalfont\textbf{Mensaje} & \textbf{Significado} & \textbf{Qué hace MATLAB} & \textbf{Voz} \\
\midrule
\endhead
M + modo & Modo de la tarjeta & Cambia de panel; en práctica envía \cod{L} & ``Modo lectura'' \\
D + letra & Letra mostrada & Letra gigante en el panel de lectura & La letra \\
E & Fin de palabra & ``Fin de la palabra. Escribe otra.'' & --- \\
K + letra & Letra tecleada & \cod{letraEscrita}: agrega al texto & Letra o palabra \\
T + letra & Letra objetivo & Muestra la letra y ``Letra 3 de 10'' & ``Letra m'' \\
O c hh mm & Acierto & Verde, marcador, gráfica & ``Correcto'' \\
X c hh mm & Error & Rojo, marcador, gráfica & ``Incorrecto'' \\
F hh mm & Fin de sesión & Resultado final con precisión & Resumen \\
L hh & Letras confirmadas & Lo anota en el registro & --- \\
S hh mm & Marcador & Lo anota en el registro & --- \\
!L !U !M & Error & Mensaje de error en el registro & --- \\
\bottomrule
\end{longtable}

\subsection{Caso \texttt{'M'}: cambio de modo}
\begin{lstlisting}[firstnumber=198]
            case 'M'                                     % modo (switches)
                mostrarModo(c2);
                registrar(['La tarjeta esta en modo ' nombreModo(c2)]);
                hablar(['Modo ' lower(nombreModo(c2))]);
                if c2 == 'P'
                    enviarN();          % sesion nueva: decirle cuantas letras
                end
\end{lstlisting}
Ejemplo: llega \cod{MP} $\rightarrow$ muestra el panel de práctica, dice ``modo
práctica'' y envía \cod{L10} para que la FPGA sepa cuántas letras usar.

\subsection{Caso \texttt{'D'}: letra mostrada (lectura)}
\begin{lstlisting}[firstnumber=206]
            case 'D'                                     % lectura
                lblLetra.Text = upper(c2);
                lblProg.Text  = 'Avanza con KEY1 (NEXT) en la tarjeta.';
                hablar(c2);
\end{lstlisting}
\cod{upper} la muestra en mayúscula. Ejemplo: \cod{Dh} $\rightarrow$ aparece una
\textbf{H} grande y la voz dice ``h''.

\subsection{Caso \texttt{'E'}: fin de palabra}
\begin{lstlisting}[firstnumber=211]
            case 'E'
                lblLetra.Text = '-';
                lblProg.Text  = 'Fin de la palabra. Escribe otra.';
                registrar('Fin de la palabra.');
\end{lstlisting}

\subsection{Caso \texttt{'K'}: letra tecleada (escritura)}
\begin{lstlisting}[firstnumber=216]
            case 'K'                                     % escritura
                letraEscrita(c2);
\end{lstlisting}
Delega todo a \cod{letraEscrita} (sección~\ref{sec:escritura}).

\subsection{Caso \texttt{'T'}: nueva letra objetivo (práctica)}
\begin{lstlisting}[firstnumber=219]
            case 'T'                                     % practica: letra nueva
                lblObj.Text = upper(c2);
                lblObj.FontColor = [0 0 0];
                if nLetras > 0
                    lblRes.Text = sprintf('Letra %d de %d: escribela en el teclado', ...
                                          intentos + 1, nLetras);
                else
                    lblRes.Text = 'Escribe esta letra en el teclado Braille';
                end
                hablar(['Letra ' c2]);
\end{lstlisting}
\begin{itemize}
  \item La letra se pinta en negro (vuelve al color normal después de un acierto o error).
  \item \cod{sprintf} arma un texto con números: cada \cod{\%d} se reemplaza por un entero.
        Con \cod{intentos = 2} y \cod{nLetras = 10}: ``Letra 3 de 10''.
  \item Si \cod{nLetras = 0} (sin límite), no muestra el contador.
\end{itemize}

\subsection{Caso \texttt{'O'} y \texttt{'X'}: resultado del intento}
\begin{lstlisting}[firstnumber=230]
            case {'O', 'X'}                              % practica: resultado
                ok = (m(1) == 'O');
                h  = hex2dec(char(m(3:4)));
                e  = hex2dec(char(m(5:6)));
                acc = 100 * h / max(1, h + e);
                intentos = h + e;
                histAcc(end + 1) = acc;
                if ok
                    lblObj.FontColor = [0 0.6 0];
                    lblRes.Text = 'Correcto!';
                    hablar('Correcto');
                else
                    lblObj.FontColor = [0.8 0 0];
                    lblRes.Text = sprintf('Incorrecto: escribiste %s', upper(c2));
                    hablar('Incorrecto');
                end
                lblStats.Text = sprintf('Aciertos: %d   Errores: %d   Precision: %.0f %%', h, e, acc);
                plot(ax, 1:numel(histAcc), histAcc, '-o', 'LineWidth', 1.5);
                ylim(ax, [0 100]); grid(ax, 'on');
                registrar(sprintf('%s  (aciertos %d, errores %d)', lblRes.Text, h, e));
\end{lstlisting}

\begin{longtable}{>{\ttfamily}p{4.4cm} p{9.7cm}}
\toprule
\normalfont\textbf{Código} & \textbf{Explicación} \\
\midrule
\endhead
case \{'O', 'X'\} & Un mismo bloque atiende los dos casos. \\
ok = (m(1) == 'O') & Verdadero si fue acierto. \\
hex2dec(char(m(3:4))) & Convierte los dígitos hexadecimales en número: \cod{'03'} $\rightarrow$ 3, \cod{'0C'} $\rightarrow$ 12. Es lo contrario de \cod{hexc} en la FPGA. \\
100 * h / max(1, h + e) & Precisión en porcentaje. \cod{max(1, ...)} evita dividir entre cero. \\
histAcc(end + 1) = acc & Agrega la precisión al final de la historia. \\
FontColor = [0 0.6 0] & Verde si acierto; \cod{[0.8 0 0]} rojo si error. \\
\%.0f \%\% & En \cod{sprintf}: número sin decimales y el símbolo \%. \\
plot(ax, x, y, '-o') & Redibuja la gráfica: eje X = intento 1, 2, 3\ldots; eje Y = precisión; línea con círculos. \\
ylim(ax, [0 100]) & Fija el eje Y de 0 a 100 \%. \\
\bottomrule
\end{longtable}

\begin{idea}[Reparto de trabajo]
La FPGA \textbf{cuenta} aciertos y errores (es la que sabe si acertaste). MATLAB
\textbf{calcula el porcentaje} y \textbf{grafica}, porque dividir con decimales en
hardware es costoso y en la PC es una línea de código.
\end{idea}

\subsection{Caso \texttt{'F'}: fin de sesión}
\begin{lstlisting}[firstnumber=251]
            case 'F'                                     % practica: fin de sesion
                h   = hex2dec(char(m(2:3)));
                e   = hex2dec(char(m(4:5)));
                acc = 100 * h / max(1, h + e);
                lblObj.Text = char(10003);
                lblObj.FontColor = [0 0.4 0.8];
                lblRes.Text = 'Sesion terminada.  KEY1 (NEXT) = otra sesion';
                lblStats.Text = sprintf('FINAL  -  Aciertos: %d   Errores: %d   Precision: %.0f %%', h, e, acc);
                registrar(sprintf('Sesion terminada: aciertos %d, errores %d, precision %.0f %%', h, e, acc));
                hablar(sprintf('Practica terminada. %d aciertos, %d errores. Precision %.0f por ciento', h, e, acc));
\end{lstlisting}
\begin{ojo}
Aquí los números están en las posiciones \textbf{2--3 y 4--5}, no en 3--4 y 5--6,
porque el mensaje \cod{F} no lleva letra: \cod{F0802} = 8 aciertos y 2 errores.
\cod{char(10003)} es el símbolo $\checkmark$.
\end{ojo}

\subsection{Caso \texttt{'L'}: confirmación de letras por sesión}
\begin{lstlisting}[firstnumber=262]
            case 'L'                                     % letras por sesion
                nLetras = hex2dec(char(m(2:3)));
                if nLetras == 0
                    registrar('Practica sin limite de letras.');
                else
                    registrar(sprintf('La FPGA hara %d letras por sesion.', nLetras));
                end
\end{lstlisting}
Ejemplo: \cod{L0A} $\rightarrow$ \cod{nLetras = 10}.

\subsection{Caso \texttt{'S'}: marcador}
\begin{lstlisting}[firstnumber=270]
            case 'S'
                if numel(m) >= 5
                    registrar(sprintf('Marcador FPGA: aciertos %d, errores %d', ...
                        hex2dec(char(m(2:3))), hex2dec(char(m(4:5)))));
                end
\end{lstlisting}
Revisa que el mensaje tenga al menos 5 caracteres antes de leer las posiciones,
para no salirse del arreglo si llegara incompleto.

\subsection{Caso \texttt{'!'} y \texttt{otherwise}: errores y mensajes desconocidos}
\begin{lstlisting}[firstnumber=276]
            case '!'
                switch c2
                    case 'L', registrar('Error: palabra demasiado larga.');
                    case 'M', registrar('Pon la tarjeta en modo lectura (SW1 SW0 = 01).');
                    otherwise, registrar(['Error: ' char(m)]);
                end

            otherwise
                registrar(['Mensaje: ' char(m)]);
\end{lstlisting}
\cod{otherwise} anota cualquier mensaje no reconocido tal cual. Es útil para depurar.

% =====================================================================
\section{Cambio de panel}

\subsection{\texttt{mostrarModo}}
\linea{289--306}
\begin{lstlisting}[firstnumber=289]
    function mostrarModo(m)
        modoAct = m;
        pI.Visible = (m == 'I');
        pR.Visible = (m == 'R');
        pW.Visible = (m == 'W');
        pP.Visible = (m == 'P');
        lblModo.Text = ['MODO: ' nombreModo(m)];
        switch m
            case 'R'
                lblLetra.Text = '-';
            case 'P'                                    % la FPGA borra el marcador
                histAcc = [];
                intentos = 0;
                cla(ax); ylim(ax, [0 100]); grid(ax, 'on');
                lblObj.Text = '-';
                lblStats.Text = 'Aciertos: 0   Errores: 0   Precision: 0 %';
        end
    end
\end{lstlisting}
\begin{itemize}
  \item \cod{(m == 'I')} vale verdadero o falso: solo un panel queda visible.
  \item Al entrar a práctica reinicia la historia, los intentos y la gráfica
        (\cod{cla} = \emph{clear axes}), igual que la FPGA reinicia su marcador.
\end{itemize}

\subsection{\texttt{nombreModo}}
\linea{308--315}
Convierte la letra en un nombre legible: \cod{'R'} $\rightarrow$ \cod{'LECTURA'},
\cod{'W'} $\rightarrow$ \cod{'ESCRITURA'}, \cod{'P'} $\rightarrow$ \cod{'PRACTICA'},
cualquier otro $\rightarrow$ \cod{'REPOSO'}.

% =====================================================================
\section{Lo que MATLAB envía}

\subsection{Práctica: \texttt{enviarN} y \texttt{cambiarN}}
\linea{321--332}
\begin{lstlisting}[firstnumber=321]
    function enviarN()
        nLetras = spnN.Value;
        enviar("L" + string(spnN.Value));
    end

    function cambiarN(~, ~)
        if ~isempty(s) && modoAct == 'P'
            enviarN();
        end
    end
\end{lstlisting}

\begin{center}
\begin{tabular}{>{\raggedright}p{6.5cm} >{\raggedright\arraybackslash}p{6.5cm}}
\toprule
\textbf{Cuándo se envía \cod{L}} & \textbf{Quién lo dispara} \\
\midrule
Al entrar al modo práctica & \cod{case 'M'} con \cod{c2 == 'P'} \\
Al cambiar el número durante la práctica & \cod{cambiarN} (callback del selector) \\
\bottomrule
\end{tabular}
\end{center}
Con strings, \cod{"L" + string(10)} produce \cod{"L10"}. \cod{cambiarN} solo envía
si hay conexión \textbf{y} la tarjeta está en práctica.

\subsection{Lectura: \texttt{enviarPalabra}}
\linea{338--348}
\begin{lstlisting}[firstnumber=338]
    function enviarPalabra(~, ~)
        w = strtrim(edPal.Value);
        if isempty(w), return; end
        if numel(w) > MAX_LEN
            uialert(fig, sprintf('Maximo %d caracteres.', MAX_LEN), 'Palabra muy larga');
            return;
        end
        registrar(['Enviando palabra: ' w]);
        hablar(w);
        enviar("R" + string(w));
    end
\end{lstlisting}
\begin{enumerate}
  \item \cod{strtrim} quita espacios al inicio y al final.
  \item Si está vacía, no hace nada.
  \item Si tiene más de 16 caracteres, avisa con una ventana y \textbf{no la envía}.
  \item La anota, la dice en voz alta y envía \cod{R} + palabra.
\end{enumerate}
\begin{idea}[Doble protección]
MATLAB revisa el largo antes de enviar para dar un aviso claro. La FPGA también lo
revisa y responde \cod{!L} por si la palabra llegara desde otro programa.
\end{idea}

% =====================================================================
\section{Escritura: \texttt{letraEscrita} y \texttt{borrarTexto}}\label{sec:escritura}
\linea{354--376}

\begin{lstlisting}[firstnumber=354]
    function letraEscrita(c)
        if c == '?'
            beep;
            lblUlt.Text = 'Patron no reconocido';
            return;
        end
        if c == ' '
            if ~isempty(palabraW), hablar(palabraW); end
            palabraW = '';
            lblUlt.Text = 'Espacio';
        else
            palabraW = [palabraW c];
            lblUlt.Text = ['Ultima letra: ' upper(c)];
            hablar(c);
        end
        texto = [texto c];
        taTexto.Value = {texto};
    end
\end{lstlisting}

\begin{center}
\begin{tabular}{l l >{\raggedright\arraybackslash}p{9.5cm}}
\toprule
\textbf{Caso} & \textbf{Llega} & \textbf{Qué hace} \\
\midrule
Patrón inválido & \cod{K?} & Suena un \cod{beep}, avisa y \textbf{no agrega nada} \\
Espacio (NEXT) & \cod{K\textvisiblespace} & Lee en voz alta la palabra completa y empieza una nueva \\
Letra normal & \cod{Ka} & La agrega a la palabra actual y la dice \\
(todos menos ?) & --- & La agrega al texto total y actualiza la pantalla \\
\bottomrule
\end{tabular}
\end{center}

\subsection*{Ejemplo: teclear ``hola'' y NEXT}
\begin{center}
\begin{tabular}{l l l l}
\toprule
\textbf{Llega} & \textbf{Voz} & \textbf{\cod{palabraW}} & \textbf{Pantalla} \\
\midrule
\cod{Kh} & ``h'' & h & h \\
\cod{Ko} & ``o'' & ho & ho \\
\cod{Kl} & ``l'' & hol & hol \\
\cod{Ka} & ``a'' & hola & hola \\
\cod{K\textvisiblespace} & ``hola'' & (vacía) & hola\textvisiblespace \\
\bottomrule
\end{tabular}
\end{center}

\cod{borrarTexto} (líneas 373--376) vacía \cod{texto}, \cod{palabraW} y el área de
texto. Solo afecta a la ventana; la FPGA no se entera.

% =====================================================================
\section{Utilidades}

\subsection{\texttt{registrar}: el historial}
\linea{382--386}
\begin{lstlisting}[firstnumber=382]
    function registrar(txt)
        linea = sprintf('[%s] %s', char(datetime('now', 'Format', 'HH:mm:ss')), txt);
        v = [{linea}; taLog.Value(:)];
        taLog.Value = v(1:min(end, 200));
    end
\end{lstlisting}
\begin{itemize}
  \item Agrega la hora: \cod{[12:05:33] Conectado a COM5.}
  \item Pone la línea nueva \textbf{arriba} de las anteriores.
  \item Conserva solo las últimas \textbf{200} líneas para que no crezca sin límite.
\end{itemize}

\subsection{\texttt{hablar}: la voz}
\linea{389--401}
\begin{lstlisting}[firstnumber=389]
    function hablar(txt)
        if ~chkVoz.Value || isempty(txt), return; end
        try
            if isempty(synth)
                NET.addAssembly('System.Speech');
                synth = System.Speech.Synthesis.SpeechSynthesizer;
            end
            synth.SpeakAsyncCancelAll();
            synth.SpeakAsync(char(txt));
        catch
            % Sin voz disponible
        end
    end
\end{lstlisting}
\begin{longtable}{>{\ttfamily}p{4.6cm} p{9.5cm}}
\toprule
\normalfont\textbf{Código} & \textbf{Explicación} \\
\midrule
\endhead
\textasciitilde chkVoz.Value & Si la casilla Voz está apagada, no habla. \\
NET.addAssembly & Carga la librería de voz integrada en Windows (.NET). \\
if isempty(synth) & Crea el sintetizador \textbf{una sola vez} y lo reutiliza. \\
SpeakAsyncCancelAll() & Interrumpe lo que estuviera diciendo, para que no se acumulen frases atrasadas. \\
SpeakAsync & Habla \textbf{sin congelar} la ventana (asíncrono). \\
try / catch & En Mac o Linux, o si falla la voz, simplemente no habla. \\
\bottomrule
\end{longtable}

\subsection{\texttt{cerrar}: salida limpia}
\linea{403--409}
\begin{lstlisting}[firstnumber=403]
    function cerrar(~, ~)
        if ~isempty(s)
            try, configureCallback(s, "off"); catch, end
            delete(s);
        end
        delete(fig);
    end
\end{lstlisting}
\begin{ojo}[¿Por qué es importante?]
Si el puerto no se libera, queda ocupado: al volver a abrir el programa
aparecería ``No se pudo abrir el puerto'' hasta reiniciar MATLAB.
\end{ojo}

% =====================================================================
\section{Resumen del protocolo desde MATLAB}

\subsection{Lo que MATLAB envía}
\begin{center}
\begin{tabular}{>{\ttfamily}l l l >{\raggedright\arraybackslash}p{4.2cm}}
\toprule
\normalfont\textbf{Envía} & \textbf{Ejemplo en bytes} & \textbf{Función} & \textbf{Cuándo} \\
\midrule
Q & \cod{81 10} & \cod{conectar} & Al conectar \\
R + palabra & \cod{82 104 111 108 97 10} & \cod{enviarPalabra} & Botón ``Enviar palabra'' \\
L + número & \cod{76 49 48 10} & \cod{enviarN} & Al entrar a práctica o cambiar el número \\
\bottomrule
\end{tabular}
\end{center}

\subsection{Formato de los mensajes recibidos}
\begin{center}
\begin{tabular}{>{\ttfamily}l l >{\raggedright\arraybackslash}p{7.2cm}}
\toprule
\normalfont\textbf{Mensaje} & \textbf{Bytes (con \cod{\textbackslash n})} & \textbf{Posiciones que lee MATLAB} \\
\midrule
MR & 3 & \cod{c2} = modo \\
Dh, Ka, Tm & 3 & \cod{c2} = letra \\
E & 2 & --- \\
Om0301 / Xk0302 & 7 & \cod{c2} letra, \cod{m(3:4)} aciertos, \cod{m(5:6)} errores \\
F0802 & 6 & \cod{m(2:3)} aciertos, \cod{m(4:5)} errores \\
L0A & 4 & \cod{m(2:3)} número \\
S0C03 & 6 & \cod{m(2:3)} aciertos, \cod{m(4:5)} errores \\
!L, !U, !M & 3 & \cod{c2} = tipo de error \\
\bottomrule
\end{tabular}
\end{center}

% =====================================================================
\section{Flujos completos por modo}

\subsection{Conexión}
\begin{longtable}{c p{4.2cm} p{4.6cm} p{4.4cm}}
\toprule
\textbf{\#} & \textbf{Usuario} & \textbf{MATLAB} & \textbf{FPGA} \\
\midrule
\endhead
1 & Ejecuta \cod{braille\_gui} & Crea la ventana, panel Reposo & --- \\
2 & Elige COM5, presiona Conectar & Abre el puerto, envía \cod{Q} & --- \\
3 & --- & --- & Responde \cod{MI} \\
4 & --- & \cod{alRecibir} $\rightarrow$ \cod{procesar} $\rightarrow$ panel Reposo & --- \\
\bottomrule
\end{longtable}

\subsection{Lectura}
\begin{longtable}{c p{4.2cm} p{4.6cm} p{4.4cm}}
\toprule
\textbf{\#} & \textbf{Usuario} & \textbf{MATLAB} & \textbf{FPGA} \\
\midrule
\endhead
1 & Pone SW = 01 & --- & Envía \cod{MR} \\
2 & --- & Panel Lectura, dice ``modo lectura'' & --- \\
3 & Escribe ``hola'', presiona Enviar & Envía \cod{Rhola} & Guarda la palabra, servos con la h \\
4 & --- & --- & Envía \cod{Dh} \\
5 & --- & H gigante, dice ``h'' & --- \\
6 & Presiona NEXT & --- & Servos con la o, envía \cod{Do} \\
7 & NEXT tras la última letra & ``Fin de la palabra'' & Servos abajo, envía \cod{E} \\
\bottomrule
\end{longtable}

\subsection{Escritura}
\begin{longtable}{c p{4.2cm} p{4.6cm} p{4.4cm}}
\toprule
\textbf{\#} & \textbf{Usuario} & \textbf{MATLAB} & \textbf{FPGA} \\
\midrule
\endhead
1 & Pone SW = 10 & Panel Escritura & Envía \cod{MW} \\
2 & Teclea puntos 1, 2, 5 & Agrega ``h'' y la dice & Reconoce h, envía \cod{Kh} \\
3 & Teclea un patrón inválido & Beep, ``Patrón no reconocido'' & Envía \cod{K?} \\
4 & Presiona NEXT & Lee la palabra completa & Envía \cod{K\textvisiblespace} \\
\bottomrule
\end{longtable}

\subsection{Práctica}
\begin{longtable}{c p{4.2cm} p{4.6cm} p{4.4cm}}
\toprule
\textbf{\#} & \textbf{Usuario} & \textbf{MATLAB} & \textbf{FPGA} \\
\midrule
\endhead
1 & Pone SW = 11 & --- & Envía \cod{MP}, borra marcador \\
2 & --- & Panel Práctica, reinicia gráfica, envía \cod{L10} & Responde \cod{L0A} \\
3 & --- & --- & Elige letra al azar, envía \cod{Tm} \\
4 & --- & Muestra M, dice ``letra m'' & --- \\
5 & Teclea la m correctamente & --- & Compara, \cod{hits + 1}, envía \cod{Om0100} \\
6 & --- & M en verde, ``Correcto'', gráfica en 100 \% & Muestra la respuesta 2 s \\
7 & (repite hasta 10 letras) & --- & Envía \cod{F0802} \\
8 & --- & $\checkmark$ azul, resultado final, lo dice & Espera NEXT para otra sesión \\
\bottomrule
\end{longtable}

% =====================================================================
\section{Preguntas frecuentes (para la presentación)}

\begin{pregunta}{¿Por qué casi toda la lógica está en la FPGA y no en MATLAB?}
Porque el objetivo del proyecto es el hardware: la tarjeta funciona sola (con los
switches, botones, servos y displays). MATLAB solo agrega pantalla, voz y gráficas.
Además, así la FPGA siempre es la ``fuente de verdad'': el marcador y el modo no
pueden quedar desincronizados.
\end{pregunta}

\begin{pregunta}{¿Cómo sabe MATLAB en qué modo está la tarjeta?}
No puede ver los switches. Al conectarse envía \cod{Q} y la FPGA responde con
\cod{M} + modo. Después, cada vez que se mueven los switches, la FPGA envía
\cod{M} por su cuenta.
\end{pregunta}

\begin{pregunta}{¿Cómo sabe MATLAB dónde termina un mensaje?}
Todos los mensajes terminan con el byte 10 (\cod{\textbackslash n}).
\cod{configureCallback(s, "terminator", @alRecibir)} ejecuta \cod{alRecibir} cada
vez que llega uno, y esa función corta el buffer en cada byte 10.
\end{pregunta}

\begin{pregunta}{¿Qué pasa si dos mensajes llegan juntos o uno llega a medias?}
\cod{alRecibir} usa un buffer propio (\cod{rxbuf}) y un ciclo \cod{while}: procesa
todos los mensajes completos que haya y deja guardado el pedazo incompleto hasta
que llegue su \cod{\textbackslash n}.
\end{pregunta}

\begin{pregunta}{¿Por qué los contadores llegan en hexadecimal?}
Si se enviara el número como un byte crudo, el valor 10 se confundiría con el
\cod{\textbackslash n} y cortaría el mensaje. En hexadecimal cada dígito es 0--9 o
A--F, que nunca valen 10. MATLAB los convierte con \cod{hex2dec}.
\end{pregunta}

\begin{pregunta}{¿Dónde se calcula la precisión?}
En MATLAB: \cod{acc = 100 * h / max(1, h + e)}. La FPGA solo cuenta aciertos y
errores; dividir con decimales en hardware sería costoso.
\end{pregunta}

\begin{pregunta}{¿Cómo cambia de panel la ventana?}
Los cuatro paneles están en la misma celda de la cuadrícula. \cod{mostrarModo}
deja visible solo el del modo actual con \cod{pX.Visible = (m == 'X')}.
\end{pregunta}

\begin{pregunta}{¿Por qué las funciones pueden usar variables como \texttt{s} sin recibirlas?}
Porque son \textbf{funciones anidadas} dentro de \cod{braille\_gui}: comparten
todas sus variables.
\end{pregunta}

\begin{pregunta}{¿Cómo funciona la voz y por qué no congela la ventana?}
Usa el sintetizador de Windows (\cod{System.Speech}) mediante .NET. Con
\cod{SpeakAsync} habla en segundo plano y la ventana sigue respondiendo.
\cod{SpeakAsyncCancelAll} evita que se acumulen frases atrasadas.
\end{pregunta}

\begin{pregunta}{¿Qué pasa si el puerto no existe o está ocupado?}
\cod{serialport} está dentro de un \cod{try/catch}: en lugar de detener el
programa, se muestra una ventana de error con \cod{uialert}.
\end{pregunta}

\begin{pregunta}{¿Por qué la ñ funciona?}
\cod{uint8(char(txt))} convierte cada carácter en un solo byte; la ñ queda como 241
(Latin-1), que es el valor que la FPGA tiene en su tabla Braille.
\end{pregunta}

\begin{pregunta}{¿Por qué se valida el largo de la palabra en MATLAB si la FPGA también lo hace?}
Es doble protección: MATLAB da un aviso claro antes de enviar, y la FPGA responde
\cod{!L} si una palabra larga llegara desde otro programa.
\end{pregunta}

\begin{pregunta}{¿Qué hace el programa al cerrarse?}
\cod{cerrar} apaga el callback y libera el puerto con \cod{delete(s)}. Si no, el
puerto quedaría ocupado para la siguiente ejecución.
\end{pregunta}

% =====================================================================
\appendix
\section{Índice rápido de funciones}
\begin{longtable}{>{\ttfamily}p{3cm} p{1.8cm} p{4cm} p{4.2cm}}
\toprule
\normalfont\textbf{Función} & \textbf{Líneas} & \textbf{Se ejecuta cuando} & \textbf{Hace} \\
\midrule
\endhead
listaPuertos & 135--138 & Inicio y botón refrescar & Lista los COM disponibles \\
conectar & 140--164 & Botón Conectar & Abre o cierra el puerto, envía \cod{Q} \\
enviar & 166--172 & Lo llaman otras funciones & Envía texto + \cod{\textbackslash n} \\
alRecibir & 175--186 & Llega un \cod{\textbackslash n} & Corta mensajes y llama a \cod{procesar} \\
procesar & 192--286 & Por cada mensaje & Actualiza la ventana según el tipo \\
mostrarModo & 289--306 & Mensaje \cod{M} e inicio & Cambia el panel visible \\
nombreModo & 308--315 & Lo llaman otras funciones & Letra $\rightarrow$ nombre del modo \\
enviarN & 321--324 & Entrar a práctica & Envía \cod{L} + número \\
cambiarN & 328--332 & Cambiar el selector & Envía \cod{L} si está en práctica \\
enviarPalabra & 338--348 & Botón Enviar palabra & Valida y envía \cod{R} + palabra \\
letraEscrita & 354--371 & Mensaje \cod{K} & Agrega la letra al texto \\
borrarTexto & 373--376 & Botón Borrar & Limpia el texto \\
registrar & 382--386 & Lo llaman otras funciones & Anota en el historial con hora \\
hablar & 389--401 & Lo llaman otras funciones & Voz asíncrona de Windows \\
cerrar & 403--409 & Cerrar la ventana & Libera el puerto \\
\bottomrule
\end{longtable}

\end{document}
